% Generated draft: requires mathematical and source review.
\documentclass[11pt]{article}
\usepackage[margin=25mm]{geometry}
\usepackage{fontspec}
\setmainfont{Latin Modern Roman}
\usepackage{amsmath,amssymb,mathtools,unicode-math}
\setmathfont{Latin Modern Math}
\usepackage{enumitem,needspace,xurl,hyperref}
\setlength{\parindent}{0pt}
\setlength{\parskip}{5pt}
\setlength{\emergencystretch}{3em}
\newcommand{\N}{\mathbb{N}}
\newcommand{\Z}{\mathbb{Z}}
\newcommand{\Q}{\mathbb{Q}}
\newcommand{\R}{\mathbb{R}}
\newcommand{\C}{\mathbb{C}}
\newcommand{\F}{\mathbb{F}}
\newcommand{\A}{\mathbb{A}}
\newcommand{\PP}{\mathbb{P}}
\newcommand{\E}{\mathbb{E}}
\newcommand{\eps}{\varepsilon}
\newcommand{\dd}{\,\mathrm{d}}
\begin{document}
\section*{UnsolvedMath ID: 700390}
\textbf{Source problem code:} catalogue entry
\subsection{Definitions and mathematical statement}
Let \(G=(X,\cdot,{}^{-1},1_G)\) be an arbitrary group, and let \(v\in F(x,y)\) be a group word in two variables, where \(F(x,y)\) is the free group on \(x,y\). Evaluate \(v\) using the original group operations of \(G\), and define
\[
a\odot b=v^G(a,b)\qquad(a,b\in X).
\]
Assume that \(\odot\) makes \(X\) into a group, denoted by \(G_v\). This assumption concerns the chosen group \(G\) and word \(v\); it does not require \(v\) to define a group operation on every group.

For any group \(H\), a group identity \(w=1\), with \(w\in F(x_1,\ldots,x_n)\), holds in \(H\) if \(w^H(h_1,\ldots,h_n)=1_H\) for every \((h_1,\ldots,h_n)\in H^n\). Define the variety generated by \(G\) by
\[
\operatorname{Var}(G)=\{H:\ H\text{ is a group satisfying every group identity holding in }G\}.
\]
Is it true, for every such \(G\) and \(v\), that
\[
G_v\in\operatorname{Var}(G)?
\]
Equivalently, must the implication
\[
\left(\forall (g_1,\ldots,g_n)\in X^n,\quad w^G(g_1,\ldots,g_n)=1_G\right)
\Longrightarrow
\left(\forall (g_1,\ldots,g_n)\in X^n,\quad w^{G_v}(g_1,\ldots,g_n)=1_{G_v}\right)
\]
hold for every integer \(n\geq 1\) and every group word \(w\in F(x_1,\ldots,x_n)\)? Here evaluation in \(G_v\) uses its induced multiplication, inverse, and identity.

The supplied record labels this problem open; that status is not independently verified here. Its attribution is unresolved: the statement names C. D. H. Cooper, while the proposer metadata and background name E. I. Khukhro.
\subsection{Short English statement}
If a word in the original operations of a group defines another group operation on the same set, must the resulting group satisfy every group identity satisfied by the original group?
\subsection{Sources}
\begin{itemize}
\item Local UnsolvedMath registry, permanent ID \(700390\), exact source code \texttt{catalogue entry}, title ``Kourovka Notebook Problem 6.47.'' The supplied record labels the problem open and gives the proposed year as 1978.
\item Supplied source reference: \emph{The Kourovka Notebook (Unsolved Problems in Group Theory)}, 6th issue (1978), Problem 6.47; 21st issue edition (2026), \texttt{arXiv:1401.0300v48}. The supplied collection description names E. I. Khukhro and V. D. Mazurov as editors.
\item Supplied attribution: C. D. H. Cooper appears at the beginning of the statement; E. I. Khukhro appears in the background and proposer metadata. The record does not explain their respective roles.
\end{itemize}
\end{document}
